% --- CORRECCIÓN IMPORTANTE EN LA PRIMERA LÍNEA ---
% Agregamos "xcolor=table" para que funcione \rowcolor sin errores
\documentclass[aspectratio=169, 10pt, xcolor=table]{beamer}

% --- Configuración del Tema (estilo profesional como el ejemplo) ---
\usetheme{Madrid}
\usecolortheme{dolphin}

% --- Paquetes necesarios ---
\usepackage[utf8]{inputenc}
\usepackage[spanish]{babel}
\usepackage{amsmath, amssymb}
\usepackage{booktabs}
\usepackage{graphicx}
\usepackage{listings}
\usepackage{tikz}
\usetikzlibrary{arrows.meta, positioning, shapes.geometric}

% --- Ruta de Imágenes (crea la carpeta "imagenes" y coloca ahí tus archivos) ---
\graphicspath{{./imagenes/}}

% --- Estilo para código Python (configurado para evitar problemas con UTF-8) ---
\definecolor{codegreen}{rgb}{0,0.6,0}
\definecolor{codegray}{rgb}{0.5,0.5,0.5}
\definecolor{codepurple}{rgb}{0.58,0,0.82}
\definecolor{backcolour}{rgb}{0.95,0.95,0.92}
\lstset{
    language=Python,
    backgroundcolor=\color{backcolour},
    commentstyle=\color{codegreen},
    keywordstyle=\color{blue},
    numberstyle=\tiny\color{codegray},
    stringstyle=\color{codepurple},
    basicstyle=\ttfamily\scriptsize,
    breaklines=true,
    frame=single,
    showstringspaces=false,
    literate={á}{{\'a}}1 {é}{{\'e}}1 {í}{{\'i}}1 {ó}{{\'o}}1 {ú}{{\'u}}1 {ñ}{{\~n}}1 {ü}{{\"u}}1
}

% --- Pie de página personalizado ---
\makeatletter

\setbeamertemplate{footline}{
  \leavevmode%
  \hbox{%
  \begin{beamercolorbox}[wd=.333333\paperwidth,ht=2.25ex,dp=1ex,center]{author in head/foot}%
    \usebeamerfont{author in head/foot}\insertshortauthor
  \end{beamercolorbox}%
  \begin{beamercolorbox}[wd=.333333\paperwidth,ht=2.25ex,dp=1ex,center]{title in head/foot}%
    \usebeamerfont{title in head/foot}Sesi\'on 4
  \end{beamercolorbox}%
  \begin{beamercolorbox}[wd=.333333\paperwidth,ht=2.25ex,dp=1ex,right]{date in head/foot}%
    \usebeamerfont{date in head/foot}CÁLCULO Y ÁLGEBRA LINEAL \hfill \insertframenumber/\inserttotalframenumber \hspace*{0.3cm}
  \end{beamercolorbox}}%
  \vskip0pt%
}

\makeatother

% --- Datos de la portada ---
\title[SESI\'ON 4 CÁLCULO]{\textbf{Sesi\'on 4}}
\subtitle{Sistemas de Ecuaciones y la Base de la Regresión}
\author[Prof. C. S\'anchez]{Carlos C\'esar S\'anchez Coronel}
\institute{\textbf{CÁLCULO Y ÁLGEBRA LINEAL PARA IA}}
\date{}

% Logo opcional (descomentar si tienes la imagen)
% \titlegraphic{\includegraphics[height=1.5cm]{logo.png}}

% --- Eliminar la barra de navegación (opcional) ---
\setbeamertemplate{navigation symbols}{}

% --- Índice al inicio de cada sección ---
\AtBeginSection[]{
  \begin{frame}
    \frametitle{Contenido}
    \tableofcontents[currentsection]
  \end{frame}
}

% ============================================================================
% INICIO DEL DOCUMENTO
% ============================================================================
\begin{document}

% --- Portada ---
\begin{frame}
    \titlepage
\end{frame}

% --- Índice general ---
\begin{frame}{Contenido}
    \tableofcontents
\end{frame}

% ============================================================================
% SECCIÓN 1: REPRESENTACIÓN MATRICIAL DE SISTEMAS LINEALES
% ============================================================================
\section{Representación Matricial de Sistemas Lineales}

\begin{frame}{Sistemas de ecuaciones lineales}
    \begin{block}{Forma general}
        \[
        \begin{cases}
            a_{11}x_1 + a_{12}x_2 + \cdots + a_{1n}x_n = b_1 \\
            a_{21}x_1 + a_{22}x_2 + \cdots + a_{2n}x_n = b_2 \\
            \qquad \vdots \\
            a_{m1}x_1 + a_{m2}x_2 + \cdots + a_{mn}x_n = b_m
        \end{cases}
        \]
    \end{block}
    \begin{block}{Forma matricial}
        \[
        \mathbf{A}\mathbf{x} = \mathbf{b},\quad \mathbf{A}\in\mathbb{R}^{m\times n},\ \mathbf{x}\in\mathbb{R}^n,\ \mathbf{b}\in\mathbb{R}^m
        \]
    \end{block}
\end{frame}

\begin{frame}{Interpretación geométrica}
    \begin{itemize}
        \item Cada ecuación representa un hiperplano en $\mathbb{R}^n$.
        \item Para $n=2$, cada ecuación es una recta en el plano.
        \item La solución es la intersección de todos los hiperplanos.
    \end{itemize}
    \centering
    \begin{tikzpicture}[scale=0.8]
        \draw[->] (-0.5,0) -- (4,0) node[right] {$x$};
        \draw[->] (0,-0.5) -- (0,4) node[above] {$y$};
        \draw[thick, blue] (0.5,3.5) -- (3.5,0.5) node[right] {$x+y=4$};
        \draw[thick, red] (0.5,1) -- (3.5,3) node[right] {$x-y=2$};
        \fill[black] (3,1) circle (2pt) node[above right] {$(3,1)$};
    \end{tikzpicture}
\end{frame}

% ============================================================================
% SECCIÓN 2: SISTEMAS CUADRADOS (m = n)
% ============================================================================
\section{Sistemas Cuadrados (m = n)}

\begin{frame}{Eliminación Gaussiana}
    \begin{itemize}
        \item Transformar el sistema en uno triangular superior mediante operaciones elementales.
        \item Luego resolver con sustitución regresiva.
    \end{itemize}
    \begin{exampleblock}{Ejemplo}
        \[
        \begin{cases}
            2x + y - z = 8 \\
            -3x - y + 2z = -11 \\
            -2x + y + 2z = -3
        \end{cases}
        \]
        Matriz aumentada:
        \[
        \left[\begin{array}{ccc|c}
            2 & 1 & -1 & 8 \\
            -3 & -1 & 2 & -11 \\
            -2 & 1 & 2 & -3
        \end{array}\right]
        \]
    \end{exampleblock}
\end{frame}

\begin{frame}{Eliminación paso a paso}
    \begin{enumerate}
        \item Pivote $a_{11}=2$. Hacer ceros debajo:\\
        $F_2 \leftarrow F_2 + \frac{3}{2}F_1$, $F_3 \leftarrow F_3 + F_1$.
        \[
        \left[\begin{array}{ccc|c}
            2 & 1 & -1 & 8 \\
            0 & 0.5 & 0.5 & 1 \\
            0 & 2 & 1 & 5
        \end{array}\right]
        \]
        \item Intercambiar filas 2 y 3.
        \[
        \left[\begin{array}{ccc|c}
            2 & 1 & -1 & 8 \\
            0 & 2 & 1 & 5 \\
            0 & 0.5 & 0.5 & 1
        \end{array}\right]
        \]
    \end{enumerate}
\end{frame}

\begin{frame}{Eliminación y sustitución regresiva}
    \begin{enumerate}
        \setcounter{enumi}{2}
        \item Hacer cero en $(3,2)$: $F_3 \leftarrow F_3 - 0.25F_2$.
        \[
        \left[\begin{array}{ccc|c}
            2 & 1 & -1 & 8 \\
            0 & 2 & 1 & 5 \\
            0 & 0 & 0.25 & -0.25
        \end{array}\right]
        \]
        \item Resolver de abajo arriba:
        \begin{align*}
            0.25z &= -0.25 \Rightarrow z = -1 \\
            2y + 1(-1) &= 5 \Rightarrow y = 3 \\
            2x + 3 - (-1) &= 8 \Rightarrow x = 2
        \end{align*}
        Solución: $\mathbf{x} = (2,3,-1)$.
    \end{enumerate}
\end{frame}

\begin{frame}[fragile]{Resolución con NumPy}
    \begin{lstlisting}
import numpy as np

A = np.array([[2,1,-1],[-3,-1,2],[-2,1,2]], dtype=float)
b = np.array([8,-11,-3])
x = np.linalg.solve(A, b)
print("Solución:", x)          # [ 2.  3. -1.]
# Verificar
print("A @ x - b:", np.dot(A,x)-b)  # casi cero
    \end{lstlisting}
\end{frame}

% ============================================================================
% SECCIÓN 3: SISTEMAS SOBREDETERMINADOS (m > n) Y MÍNIMOS CUADRADOS
% ============================================================================
\section{Sistemas Sobredeterminados (m > n)}

\begin{frame}{Problema de mínimos cuadrados}
    \begin{itemize}
        \item Cuando $m > n$, normalmente no hay solución exacta.
        \item Buscamos $\mathbf{x}$ que minimice el error cuadrático:
        \[
        \|\mathbf{A}\mathbf{x} - \mathbf{b}\|_2^2 = \sum_{i=1}^m \left( \sum_{j=1}^n A_{ij}x_j - b_i \right)^2
        \]
        \item Solución: resolver las \textbf{ecuaciones normales}
        \[
        \mathbf{A}^\top \mathbf{A} \mathbf{x} = \mathbf{A}^\top \mathbf{b}
        \]
        \[
        \mathbf{x} = (\mathbf{A}^\top \mathbf{A})^{-1}\mathbf{A}^\top \mathbf{b}
        \]
    \end{itemize}
\end{frame}

\begin{frame}{Interpretación geométrica}
    \begin{itemize}
        \item La solución proyecta $\mathbf{b}$ sobre el espacio columna de $\mathbf{A}$.
        \item El residuo $\mathbf{r} = \mathbf{b} - \mathbf{A}\mathbf{x}$ es ortogonal a dicho espacio.
    \end{itemize}
    \centering
    \begin{tikzpicture}[scale=1.2]
        \draw[->] (0,0) -- (3,0) node[right] {$col_1$};
        \draw[->] (0,0) -- (0,3) node[above] {$col_2$};
        \draw[thick, blue] (0,0) -- (2.5,2) node[above right] {Espacio columna};
        \fill[red] (2.5,2.5) circle (2pt) node[above] {$\mathbf{b}$};
        \fill[blue] (2,1.6) circle (2pt) node[right] {$\mathbf{A}\mathbf{x}$};
        \draw[dashed] (2.5,2.5) -- (2,1.6);
    \end{tikzpicture}
\end{frame}

\begin{frame}{Ejemplo: Ajuste de una recta por mínimos cuadrados}
    \begin{itemize}
        \item Modelo: $y = \beta_0 + \beta_1 x$.
        \item Para $m$ puntos $(x_i, y_i)$, el sistema es:
        \[
        \begin{pmatrix} 1 & x_1 \\ 1 & x_2 \\ \vdots & \vdots \\ 1 & x_m \end{pmatrix}
        \begin{pmatrix} \beta_0 \\ \beta_1 \end{pmatrix}
        \approx \begin{pmatrix} y_1 \\ y_2 \\ \vdots \\ y_m \end{pmatrix}
        \]
        \item Se resuelve por mínimos cuadrados.
    \end{itemize}
\end{frame}

\begin{frame}[fragile]{Implementación con ecuaciones normales}
    \begin{lstlisting}
import numpy as np
import matplotlib.pyplot as plt

x = np.array([0,1,2,3,4,5])
y = np.array([1,3,2,5,4,6])

A = np.vstack([np.ones_like(x), x]).T  # matriz de diseño
ATA = A.T @ A
ATb = A.T @ y
beta = np.linalg.solve(ATA, ATb)
print(f"Intercepto: {beta[0]:.4f}, Pendiente: {beta[1]:.4f}")
    \end{lstlisting}
\end{frame}

\begin{frame}[fragile]{Gráfica del ajuste}
    \begin{lstlisting}
plt.scatter(x, y, label='Datos')
x_line = np.linspace(0,5,100)
y_line = beta[0] + beta[1] * x_line
plt.plot(x_line, y_line, 'r-', label='Ajuste')
plt.xlabel('x'); plt.ylabel('y'); plt.legend(); plt.grid()
plt.show()
    \end{lstlisting}
    \centering
    \includegraphics[width=0.5\textwidth]{ejemplo_recta.png} % si la imagen existe
\end{frame}

\begin{frame}[fragile]{Métodos alternativos: lstsq y sklearn}
    \begin{lstlisting}
# Con np.linalg.lstsq (usa QR, más estable)
beta_qr, _, _, _ = np.linalg.lstsq(A, y, rcond=None)

# Con scikit-learn
from sklearn.linear_model import LinearRegression
model = LinearRegression()
model.fit(x.reshape(-1,1), y)
print("Intercepto sklearn:", model.intercept_)
print("Pendiente sklearn:", model.coef_[0])
    \end{lstlisting}
\end{frame}

% ============================================================================
% SECCIÓN 4: APLICACIONES EN IA Y CURSOS FUTUROS
% ============================================================================
\section{Aplicaciones en IA y Cursos Futuros}

\begin{frame}{Regresión lineal en machine learning}
    \begin{itemize}
        \item Es el modelo más simple y fundacional.
        \item La solución de mínimos cuadrados es la base de la regresión lineal ordinaria (OLS).
        \item Cuando se añade regularización L2, se obtiene regresión \textbf{Ridge}.
    \end{itemize}
\end{frame}

\begin{frame}{Otras áreas de aplicación}
    \begin{itemize}
        \item \textbf{Robótica}: cinemática, filtros de Kalman.
        \item \textbf{Optimización}: problemas de mínimos cuadrados como caso convexo.
        \item \textbf{Procesamiento de señales}: filtros adaptativos (RLS).
        \item \textbf{Finanzas}: modelos de factores, estimación de betas.
    \end{itemize}
\end{frame}

% ============================================================================
% SECCIÓN 5: EJERCICIO INTEGRADOR
% ============================================================================
\section{Ejercicio Integrador: Comparación de Métodos}

\begin{frame}[fragile]{Comparación de métodos de mínimos cuadrados}
    \begin{lstlisting}
import numpy as np
from sklearn.linear_model import LinearRegression
import time

# Generar datos sintéticos: y = 2 + 3x + ruido
np.random.seed(42)
m = 1000
x = np.random.randn(m,1)
X = np.hstack([np.ones((m,1)), x])
beta_true = np.array([2,3])
y = X @ beta_true + np.random.randn(m,1)*0.5

# 1. Ecuaciones normales
start = time.time()
beta_norm = np.linalg.inv(X.T @ X) @ X.T @ y
t_norm = time.time() - start

# 2. lstsq (QR)
start = time.time()
beta_qr, _, _, _ = np.linalg.lstsq(X, y, rcond=None)
t_qr = time.time() - start

# 3. sklearn
start = time.time()
model = LinearRegression(fit_intercept=False)
model.fit(X, y)
beta_sk = model.coef_.flatten()
t_sk = time.time() - start

print("Verdaderos:", beta_true)
print("Ecuaciones normales:", beta_norm.flatten())
print("lstsq:", beta_qr.flatten())
print("sklearn:", beta_sk)
print("Tiempos: norm={:.6f}s, lstsq={:.6f}s, sklearn={:.6f}s".format(t_norm, t_qr, t_sk))
    \end{lstlisting}
\end{frame}

% ============================================================================
% SECCIÓN 6: RESUMEN
% ============================================================================
\section{Resumen}

\begin{frame}{Lo que hemos aprendido}
    \begin{exampleblock}{Conceptos clave}
        \begin{itemize}
            \item \textbf{Sistemas lineales}: representación matricial $\mathbf{A}\mathbf{x}=\mathbf{b}$.
            \item \textbf{Sistemas cuadrados}: eliminación gaussiana, factorización LU, $\mathbf{x}=\mathbf{A}^{-1}\mathbf{b}$.
            \item \textbf{Sistemas sobredeterminados}: mínimos cuadrados, ecuaciones normales, descomposición QR.
            \item \textbf{Regresión lineal}: ajuste de recta por mínimos cuadrados.
            \item \textbf{Implementaciones}: \texttt{np.linalg.solve}, \texttt{np.linalg.lstsq}, \texttt{sklearn.linear\_model.LinearRegression}.
        \end{itemize}
    \end{exampleblock}
\end{frame}

\begin{frame}{Conexión con futuros cursos}
    \begin{itemize}
        \item \textbf{Machine Learning}: regresión lineal, ridge, lasso.
        \item \textbf{Optimización}: problemas de mínimos cuadrados y gradiente descendente.
        \item \textbf{Señales y control}: filtros adaptativos, identificación de sistemas.
        \item \textbf{Deep Learning}: las redes lineales se entrenan con estos conceptos.
    \end{itemize}
\end{frame}

% --- Diapositiva final ---
\begin{frame}
    \centering
    \Huge \textbf{¡Gracias!}

\end{frame}

\end{document}